% Source: 32_Chapter_23_Judgment_In_the_Process_of_Time_Wyclif_.txt
\bookchapter[Judgment in the Process of Time]{CHAPTER 23}{JUDGMENT IN THE PROCESS OF TIME}{Wyclif and Hus}

\leadin{I}{N HIS INSIGHTFUL LITTLE} volume Christianity and History, Herbert Butterfield observes
that human systems rise and thrive and then fall because the processes of time have their own
built-in “judgment.” Institutions, which at first glance seem to be quite worthy, eventually crumble
to ruins because the centuries themselves bring out the flaws.

The Babylonian Captivity at Avignon and the Great Schism of the papacy that followed revealed
fundamental flaws. Basic reforms were in order, but after the failure of the conciliar movement no
significant reforms came from within the Church of Rome. The idea that the papal office was the
channel of God’s will died slowly. Men believed that the papacy was essential not only for the
religious life of men but as a means of sanctioning political rule. He was the center of the
Christian society on earth.

Two brave souls—John Wyclif, an Englishman, and John Hus, a Czech—dared to toy with the idea that
the Christian church was something other than a visible organization on earth headed by the pope.
They paid dearly for even raising the possibility, but they saw clearly that the hour had come for
judgment to fall upon the house of God. Who were they? And how did they point the way to the future?

\section{THE ENGLISH ZEALOT, WYCLIF}

John Wyclif was a zealot. Like most zealots, he despised neutrality, and he got as much as he gave.
Since his day, men have either praised him to high heaven or consigned him to outer darkness.

Confusion about the real Wyclif is understandable. We really know little about the man. He had an
exasperating habit of burying his character beneath pages of scholarly discourse, and while
historians can follow his arguments well enough, few claim they understand him.

Wyclif’s early life is as obscure as his personality. We are not even sure of his date of birth. He
was reared in northern England but only emerged from the medieval mists as a student at Oxford. He
secured his doctoral degree in 1372 and rose immediately to prominence as the leading professor at
the university.

The hottest issue at that hour was dominion or lordship over men. All thinkers agreed that lordship
arose from God. But how was this right to rule transmitted from God to earthly rulers? One widely
held view argued that lordship was only just when it derived from the Roman church. God had
entrusted the pope with universal dominion over all temporal things and persons. Any authority
exercised by sinful rulers was unlawful.

Other teachers insisted that lordship depended less on the mediation of the church than on the fact
that its possessor was in a state of grace, that is, he had committed no grievous sin. One of
Wyclif’s professors, Richard FitzRalph, had argued, “Why should the state of grace be required only
of temporal rulers? Do churchmen have the right to rule when they live in mortal sin?” If grace is
essential in a lay ruler, said FitzRalph, it is no less necessary in a churchman.

Undoubtedly influenced by his professor, Wyclif plunged into this debate and added an important
idea. He argued that the English government had the divinely assigned responsibility to correct the
abuses of the church within its realm and to relieve of office those churchmen who persisted in
their sin. The state could even seize the property of corrupt church officials.

Not surprisingly, the pope in 1377 condemned the Oxford Reformer’s teaching. The church might have
moved against Wyclif at that moment, but influential friends in England saw to it that the
condemnation never went beyond threats.

The long-range significance of Wyclif’s teaching on dominion lies in its link with the Reformation.
It was the English Reformer’s way of emphasizing the spiritual freedom of the righteous man. He is a
possessor of “a dominion founded on grace”: “God gives no lordship to His servants without first
giving Himself to them.” Every man, therefore, priest or layman, holds an equal place in the eyes of
God. This personal relation between a man and God is everything; character is the one basis of
office. The mediating priesthood and the sacrificial masses of the medieval church are no longer
essential. Thus Wyclif anticipates Luther’s doctrine of justification by faith alone. Both men
destroy the medieval barriers between God and his people.

This doctrine of “dominion founded in grace” proved to be only the first of Wyclif’s thunderbolts.
The decisive year of his reforming career was 1378, the date of the Great Schism in the papacy.
Faced with the comic-tragedy of one pope in Rome excommunicating another pope in Avignon, Wyclif
became more radical in his assessment of the church and its need for reform.

Wyclif’s early views of the papacy were shaped by his emphasis on apostolic poverty. He insisted
that those who sat in St. Peter’s chair should be, like the apostle, without silver or gold.
According to Wyclif the “Bible papacy” consisted in a poor and humble life, spent in the service of
the church, setting before God’s people an example of Christian goodness. The pope should be the
shepherd of the flock and the preacher who brings men to Christ.

Such a view leaves no room for the temporal power of the pope. The conception of the papacy as a
political force was anathema to Wyclif. He detested the trappings of power; he denounced the
worldliness and luxury of the popes.

In one sense Wyclif welcomed the Great Schism. The spectacle of two rival popes excommunicating each
other seemed to him to be a confirmation for all to see of the spiritual bankruptcy of the office
and the need to put something else in its place. As the schism continued, however, Wyclif’s view
hardened. He came to believe that the pope was Antichrist. If there were two of them railing at each
other, they simply shared the unholy title.

In a steady stream of charges Wyclif showed how far the papacy had departed from the simple faith
and practice of Christ and his disciples. “Christ is truth,” he wrote, “the pope is the principle of
falsehood. Christ lived in poverty, the pope labors for worldly magnificence. Christ refused
temporal dominion, the pope seeks it.”

The Oxford Reformer poured scorn on the idea that because Peter died in Rome therefore every bishop
of Rome is to be set above all of Christendom. By the same reasoning the Muslim might conclude that
his “prelate of Jerusalem,” where Christ died, is greater than the pope. Christ alone, said Wyclif,
is the head of the church. The papal institution is “full of poison.” It is Antichrist itself, the
man of sin who exalts himself above God. Let judgment fall!

\section{FROM REFORMER TO PROTESTANT}

Thus the Reformer took the great step of his life. He passed from an orthodox preacher eager for the
reform of the existing international Roman church, into a Protestant, to use a later term.

The Oxford scholar’s break with the papacy was part of a new idea that he had formed of the church.
He accepts the ancient division of the church into three parts: “one triumphant in heaven,” “one
militant here on earth,” and the third “asleep in purgatory.” But the church on earth he defines as
the whole number of the elect, containing “only men that shall be saved.” So absolute is his
predestinarianism that he adds that no man, not even a pope, knows whether he is of the church, or
whether he is “a limb of the fiend.” Wyclif guards his doctrine from some of its dangers by adding
that “as each man shall hope that he shall be safe in bliss, so he should suppose that he be a limb
of holy church,” and even maintains that “each man that shall be damned shall be damned by his own
guilt, and each man that is saved shall be saved by his own merit.”

From this doctrine of an invisible church of the elect, Wyclif drew some practical conclusions. The
church is a unity that knows nothing of papal primacies and hierarchies, and of the “sects” of
monks, friars, and priests; nor can the salvation of the elect be conditioned by masses,
indulgences, penance, or other devices of priestcraft.

In time Wyclif challenged the whole range of medieval beliefs and practices: pardons, indulgences,
absolutions, pilgrimages, the worship of images, the adoration of the saints, the treasury of their
merits laid up at the reserve of the pope, and the distinction between venial and mortal sins. He
retained belief in purgatory and extreme unction, though he admitted that he looked in vain in the
Bible for the institution of extreme unction. Images, he said, if they increased devotion, need not
be removed; and prayers to saints were not necessarily wrong. Confession he held to be useful,
provided it was voluntary and made to a suitable person, best of all if it were made in public.
Compulsory confession he considered “the bondage of Antichrist.” We can catch the spirit of his
revolt in his declaration that preaching “is of more value than the administration of any
sacrament.”

The standard Wyclif used to judge the Roman Church was the teachings of Scripture. “Neither the
testimony of Augustine nor Jerome,” he said, “nor any other saint should be accepted except in so
far as it was based upon Scripture.” “Christ’s law,” he held, “is best and enough, and other laws
men should not take, but as branches of God’s law.”

\figcaption{John Wyclif (1300–1384), the English reformer, denounced the worldliness of the popes
and emphasized the spiritual freedom of the righteous man.}

\figcaption{John Hus (1369–1415), the Czech reformer, viewed Christ, not the pope, as the head of
the church.}

The Oxford Reformer went even further in his assertion of the right of every man to examine the
Bible for himself: “The New Testament is of full authority, and open to the understanding of simple
men, as to the points that be most needful to salvation~.~.~.~. He that keepeth meekness and charity
hath the true understanding and perfection of all Holy Writ,” for “Christ did not write his laws on
tables, or on skins of animals, but in the hearts of men.”

In all of his puritan outcry, Wyclif aroused no hostility like that sparked by his attack upon the
traditional doctrine of transubstantiation. In the summer of 1380 he published twelve arguments
against the idea that the bread and wine of Holy Communion were transformed into the physical body
and blood of Christ. He asserts that the early church held that the consecrated elements of bread
and wine were efficacious symbols of Christ’s body and blood. Hence, Christ is present in the
elements sacramentally, not materially. The end of the sacrament is the presence of Christ in the
soul.

Wyclif’s denial of transubstantiation gave his enemies their opportunity. His support dwindled to a
small minority at Oxford. First the chancellor and a small council condemned his doctrines and
forbade him to lecture. Then the Archbishop of Canterbury, William Courtenay, followed with another
council that condemned ten of Wyclif’s doctrines as heretical. By 1382 the Reformer was effectively
silenced at Oxford.

Before his defeat at the university, however, Wyclif had turned for support to the people in
cottages and towns. His mission called for the Bible in the language of craftsmen and peasants, so
he led a handful of scholars at Oxford in the translation of the Latin Bible into the English
language and copied the methods of St. Francis and the friars.

From Oxford, as from Assisi two centuries before, Wyclif sent out “poor priests” into the byways and
village greens, sometimes even to churches, to win the souls of the neglected. Clad in russet robes
of undressed wool, without sandals, purse, or scrip, a long staff in their hand, dependent for food
and shelter on the good will of their neighbors, Wyclif’s poor priests soon became a power in the
land. Their enemies dubbed them Lollards, meaning mumblers. They carried a few pages of the
Reformer’s Bible and his tracts and sermons as they went throughout the countryside preaching the
Word of God. One panic-stricken observer claimed that “every second man” he met was a Lollard.
Despite Wyclif’s academic environment, his legacy was carried on largely by lay-proclaimers who
sought to model and proclaim a simple dependence upon the Bible. Its influence is an example of the
recurring empowerment of the Scriptures when read in the language of the people.

Wyclif gained enough support that church authorities had the good sense not to move against him. His
followers were hunted down, were expelled from Oxford, or forced to renounce their views, but
Wyclif, though driven from the university, was left to close his days in peace at his parish at
Lutterworth. He died there in 1384.

\section{THE PASSAGE TO BOHEMIA}

The movement Wyclif launched continued in England under restrictions but found an even greater
opportunity for expansion in Bohemia. The two nations were linked in 1383 by the marriage of Anne of
Bohemia and King Richard II of England, so students of both countries went back and forth between
Oxford and Prague.

The Wyclif revolt met greater success in Bohemia because it was joined to a strong national party
led by John Hus. The Czech Reformer came from peasant parents in southern Bohemia, a small town
called Husinetz. He studied theology at the University of Prague, earning both a bachelor of arts
(1394) and the master of arts (1396) before beginning his teaching in the faculty of arts and
plunging into the reform cause.

His student days introduced Hus to the philosophical writings of Wyclif, but only after his
ordination and appointment as rector and preacher at Bethlehem Chapel did he come upon Wyclif’s
religious writings. He adopted at once the English Reformer’s view of the church as an elect
company, with Christ, not the pope, its true head.

Bethlehem Chapel near the university gave Hus an unrivaled opportunity to circulate Wyclif’s
teachings, including his criticisms of the abuses of power in the papacy. On the walls were
paintings contrasting the behavior of the popes and Christ. The pope rode a horse; Christ walked
barefoot. Jesus washed the disciples’ feet; the pope preferred to have his kissed. The Chapel had
been founded in 1391 for the express purpose of encouraging the national faith in Bohemia, so Hus’s
fiery sermons in the Bohemian language fanned widespread popular support. Soon there were student
riots for and against Wyclif, much as today they might be for or against some revolutionary figure.

The Archbishop of Prague grew restless and complained to the pope about the spread of Wyclif’s
doctrines. Root out the heresy, replied the pope. So Archbishop Zbynek excommunicated Hus. As a
result a great popular tumult erupted. Hus made matters worse when he openly attacked the pope’s
sale of indulgences for support of his war against Naples. This move cost Hus the support of his
king Wenceslas, and when Prague fell under a papal interdict because of Hus, the Reformer left for
exile in southern Bohemia. During this period of retirement, Hus, drawing heavily upon Wyclif, wrote
his major work, On the Church.

The Council of Constance was now fast approaching, and Hus yielded to the urging of the emperor
Sigismund and agreed to appear at the council. He had hopes of presenting his views to the assembled
authorities, but upon his arrival he found himself instead a victim of the Inquisition.

\begin{sidebar}{Inquisition}
Inquisition is a strategy from ancient Roman law; it is a trial before a panel of judges. While one
judge may present the evidence, all the judges interrogate witnesses, decreasing the prospect that
the court might receive false testimony. Romans required confession or two eyewitnesses as evidence
in capital cases. Without witnesses or a confession the accused was sometimes tortured to secure a
confession when circumstances seemed to suggest guilt. This ancient arrangement was adopted by the
church.
\end{sidebar}

The rule of the Inquisition was simple: if sufficient witnesses testified to the guilt of the
accused, then he had to confess and renounce the errors or be burned. The reward for confession was
life imprisonment instead of the stake. In accordance with this rule, the panel of judges appointed
by the council believed the witnesses against Hus and condemned him for heresies he had never
taught.

\section{AWAY WITH HERETICS}

Hus was willing to yield himself to the teaching of the church, when instructed by Scripture in what
way his teaching was wrong. But he could not agree to recant heresies that he had always stoutly
disclaimed. For Hus truth was supreme: “I have said that I would not, for a chapel full of gold,
recede from the truth.” “I know,” he wrote in 1412, “that the truth stands and is mighty for ever,
and abides eternally, with whom there is no respect of persons.” Throughout his letters from
Constance, his chief anxiety is “lest liars should say that I have slipped back from the truth I
preached.” Few scenes in church history are more touching than Hus’s fidelity and refusal to swerve
from absolute truth, even to save his life.

For eight months he lay in prison in Constance. His letters during his last month rank among the
greatest in Christian literature. If the Reformer had added nothing to our intellectual heritage, he
would still have enriched our moral outlook.

“O most holy Christ,” he prayed,

\begin{blockquote}
draw me, weak as I am, after Thyself, for if Thou dost not draw us we cannot follow Thee. Strengthen
my spirit, that it may be willing. If the flesh is weak, let Thy grace precede us; come between and
follow, for without Thee we cannot go for Thy sake to cruel death. Give me a fearless heart, a right
faith, a firm hope, a perfect love, that for Thy sake I may lay down my life with patience and joy.
Amen.
\end{blockquote}

Finally, July 6, 1415, the day for his burning came. On the way to the place of execution he passed
through a churchyard and saw a bonfire of his books. He laughed and told the bystanders not to
believe the lies circulated about him. On arriving at the execution ground, familiarly known as the
Devil’s Place, Hus knelt and prayed. For the last time the marshal of the empire asked him if he
would recant and save his life. Said Hus, “God is my witness that the evidence against me is false.
I have never thought nor preached except with the one intention of winning men, if possible, from
their sins. In the truth of the gospel I have written, taught, and preached; today I will gladly
die.”

The Bohemian rebellion refused to die with Hus. It developed a moderate and a militant wing. The
moderates were called Utraquists, a term from Latin for “both” since their primary protest called
for freedom to receive Communion in both the bread and the cup.

The militants were called Taborites after the city in Bohemia that served as their chief stronghold.
These followers of Hus struggled against the Roman Church and the German Empire until several wars
reduced their number and influence. Despite the best efforts of the papacy to bring an end to the
Bohemian heresy, however, an independent church survived known as the Unitas Fratrum or Unity of the
Brotherhood. Until the coming of Luther, it remained a root in dry ground.

If the Church of Rome was to be reformed from within, it had ample opportunities in the fourteenth
and fifteenth centuries. But by the end of the fifteenth century, the dreams of Marsilius of Padua
had vanished, the leaders of reform by church councils were frustrated and repudiated, and the
revolts of Wyclif and Hus were crushed. The value of the period lies in the demonstration it gives
that reform of the papal church from within was impossible. The time of judgment had come. The Unity
of the Brotherhood was witness to that fact, and a promise of things to come.

\readinghead{Suggestions for Further Reading}
\begin{readinglist}
\entry *Evans, G.R. John Wyclif: Myth and Reality. Downers Grove, IL: IVP, 2006.
\entry *Fudge, Thomas A. The Magnificent Ride: The First Reformation in Hussite Bohemia, St. Andrews
  Studies in Reformation History. Aldershot, Vermont: Ashgate, 2008.
\entry Heinze, Rudolph W. Reform and Conflict: From the Medieval World to the Wars of Religion, AD
  1350–1648, Baker History of the Church. Vol. 4. Grand Rapids: Baker, 2004
\entry McFarlane, John. Wycliffe and the Beginnings of English Nonconformity. London: English
  University Press, 1952.
\entry Spinka, Matthew, ed. Advocates of Reform. Philadelphia: Westminster, 1953.
\entry ———. John Hus: A Biography. Princeton, NJ: Princeton University Press, 1968.
\entry Workman, Herbert B. The Dawn of the Reformation: The Age of Hus. London: Epworth, 1933.
\entry ———. The Dawn of the Reformation: The Age of Wyclif. London: Epworth, 1933.
\end{readinglist}
